% !PRINTED 368
\weberpage{368}
%<<CONT
least two of the $x_3, \dots x_7$ would form a pair in it. But that is not possible, since this pair would otherwise form a syzygetic triple with one of the remaining $x$.

To prove that the double tangents $x_1, x_2, x_3, x_4, x_5, x_6$ always occur in one complex, we take out two arbitrary ones of them, $x_1 x_2$, and choose another pair $y_1 y_2$ occurring in the complex $x_1 x_2$, which is possible in five ways. From these we form the syzygetic complex triple

\weq{2}{\begin{aligned}
\alpha)\ & x_1 x_2,\ y_1 y_2\\
\beta)\ & x_1 y_1,\ x_2 y_2\\
\gamma)\ & x_1 y_2,\ x_2 y_1,
\end{aligned}}

which represents five different such triples. Now in $\beta$) and $\gamma$) the $x_3, x_4, x_5, x_6, x_7$ must occur, and indeed no two of them paired. We now show first that these five $x$ cannot distribute themselves two and three over the two complexes $\beta$), $\gamma$), but only one and four. For suppose the complexes $\beta$), $\gamma$) were composed thus:

\weq{2}{\begin{aligned}
\beta)\ & x_1 y_1,\ x_2 y_2,\ x_3 y_3,\ x_4 y_4,\ x_5 y_5,\ \cdot\\
\gamma)\ & x_1 y_2,\ x_2 y_1,\ x_6 y_6,\ x_7 y_7,\ \cdot\quad\cdot\text{,}
\end{aligned}}

then we can further consider the complex

\weq{2}{\delta)\quad x_6 x_7,\ y_6 y_7\ \cdots{},}

which together with $\beta$) forms a syzygetic pair, because neither $x_1$ nor $y_1$ occurs in $\delta$), so the pair cannot be azygetic. Hence $\beta$) and $\delta$) must have a syzygetic quadruple in common, and since in two pairs of $\beta$) at least one $x$ different from $x_6, x_7$ occurs, then this $x$ must also occur in $\delta$), which is impossible, because no $x$ is syzygetic with $x_6, x_7$. There thus remains for $\beta$), $\gamma$) only one composition, like the following:

\weq{3}{\begin{aligned}
\beta)\ & x_1 y_1,\ x_2 y_2,\ x_3 y_3,\ x_4 y_4,\ x_5 y_5,\ x_6 y_6\\
\gamma)\ & x_1 y_2,\ x_2 y_1,\ x_7 y_7,\ \cdot\quad\cdot\quad\cdot
\end{aligned}}

That we make precisely this assumption, and do not take the six $x$ into $\gamma$), is no restriction, since we can, if necessary, interchange $y_1$ with $y_2$.

If now, keeping $x_1 x_2$ fixed, we let the other pairs of (2) $\alpha$) take the place of $y_1 y_2$, we obtain
%CONT>>
